\pdfbookmark[0]{Fünfter Abschnitt: Quadratische Formen und unbestimmte Gleichungen}{section-five}
\SourceLeaf{125}
\begin{center}
FÜNFTER ABSCHNITT

QUADRATISCHE FORMEN UND

UNBESTIMMTE GLEICHUNGEN ZWEITEN GRADES
\end{center}
\begin{center}\rule{.35\linewidth}{.4pt}\end{center}

\begin{center}
\textit{Ziel der Untersuchung; Definition und Bezeichnung der Formen.}
\end{center}

\Article{153}
In diesem Abschnitt werden wir vor allem Funktionen zweier Unbestimmter $x,y$ von der Gestalt
\[
axx+2bxy+cyy
\]
behandeln, wobei $a,b,c$ gegebene ganze Zahlen sind. Wir nennen sie \textit{Formen zweiten Grades} oder einfach \textit{Formen}. Auf dieser Untersuchung wird die Lösung des berühmten Problems aufbauen, alle Lösungen einer beliebigen unbestimmten Gleichung zweiten Grades mit zwei Unbekannten zu finden, ob diese nun ganzzahlige oder nur rationale Werte annehmen sollen. Der ausgezeichnete Lagrange hat dieses Problem bereits in seiner ganzen Allgemeinheit gelöst. Viele weitere Ergebnisse über die Natur der \textit{Formen} haben teils dieser große Geometer selbst, teils der ausgezeichnete Euler erstmals gefunden; andere, die schon Fermat entdeckt hatte, haben sie mit Beweisen versehen. Als wir die Formen eingehend untersuchten, boten sich uns jedoch so viele neue Ergebnisse, dass wir es der Mühe wert fanden, den ganzen Gegenstand von Anfang an neu aufzugreifen. Dazu hatten wir umso mehr Grund, als wir feststellten, dass die an vielen Stellen verstreuten Entdeckungen jener Männer nur wenigen bekannt waren; ferner, weil die Methode, nach der wir vorgehen werden, größtenteils unsere eigene ist; und schließlich,
\SourceLeaf{126}
weil unsere Ergebnisse ohne eine erneute Darstellung der ihren nicht einmal verständlich wären. Wir zweifeln nicht daran, dass auf diesem Gebiet noch viele ausgezeichnete Ergebnisse verborgen liegen, an denen andere ihre Kräfte erproben können. Die Geschichte bedeutender Ergebnisse werden wir stets an der jeweiligen Stelle darstellen.

Die Form $axx+2bxy+cyy$ werden wir, wenn es auf die Unbestimmten $x,y$ nicht ankommt, mit $(a,b,c)$ bezeichnen. Dieser Ausdruck bezeichnet also, ohne die Unbestimmten festzulegen, die Summe dreier Teile: des Produkts der gegebenen Zahl $a$ mit dem Quadrat einer beliebigen Unbestimmten; des Produkts der doppelten Zahl $b$ mit dieser und einer anderen Unbestimmten; und des Produkts der Zahl $c$ mit dem Quadrat der zweiten Unbestimmten. \textit{Zum Beispiel} bezeichnet $(1,0,2)$ die Summe eines Quadrats und eines doppelten Quadrats. Die Formen $(a,b,c)$ und $(c,b,a)$ bezeichnen zwar dasselbe, wenn wir nur auf die \textit{Teile selbst} achten; sie unterscheiden sich aber, wenn wir auch die \textit{Reihenfolge} der Teile berücksichtigen. Wir werden sie deshalb von nun an sorgfältig unterscheiden. Was wir dadurch gewinnen, wird im Folgenden deutlich genug werden.

\Needspace{8\baselineskip}
\begin{center}
\textit{Darstellung von Zahlen; Determinante.}
\end{center}

\Article{154}
Wir sagen, dass eine gegebene Zahl durch eine gegebene Form \textit{dargestellt} wird, wenn den Unbestimmten der Form solche ganzzahligen Werte zugewiesen werden, dass der Wert der Form gleich der gegebenen Zahl wird. Hier gilt folgender

\textsc{Satz.} \textit{Kann die Zahl $M$ durch die Form $(a,b,c)$ so dargestellt werden, dass die dabei verwendeten Werte der Unbestimmten teilerfremd sind, so ist $bb-ac$ ein quadratischer Rest der Zahl $M$.}

\textit{Beweis.} Die Werte der Unbestimmten seien $m,n$, also
\[
amm+2bmn+cnn=M
\]
und die Zahlen $\mu,\nu$ seien so gewählt, dass $\mu m+\nu n=1$ gilt (Artikel 40). Durch Ausmultiplizieren zeigt man leicht, dass
\[
\begin{aligned}
&amm+2bmn+cnn)(a\nu\nu-2b\mu\nu+c\mu\mu\\
&\qquad={}(\mu(mb+nc)-\nu(ma+nb))^2-(bb-ac)(m\mu+n\nu)^2
\end{aligned}
\]
oder
\[
M(a\nu\nu-2b\mu\nu+c\mu\mu)=(\mu(mb+nc)-\nu(ma+nb))^2-(bb-ac)
\]
Also gilt
\SourceLeaf{127}
\[
bb-ac\equiv(\mu(mb+nc)-\nu(ma+nb))^2\ (\mathrm{mod.}\,M)
\]
das heißt, $bb-ac$ ist ein quadratischer Rest von $M$.

Die Zahl $bb-ac$ nennen wir die \textit{Determinante} der Form $(a,b,c)$. Wie wir später zeigen werden, hängen die Eigenschaften einer Form vor allem von der Beschaffenheit dieser Zahl ab.

\begin{center}
\textit{Werte des Ausdrucks $\surd(bb-ac)(\mathit{mod.}\,M)$, zu denen eine Darstellung der Zahl $M$ durch die Form $(a,b,c)$ gehört.}
\end{center}

\Article{155}
Damit ist
\[
\mu(mb+nc)-\nu(ma+nb)
\]
ein Wert des Ausdrucks
\[
\surd(bb-ac)\ (\mathrm{mod.}\,M)
\]
Die Zahlen $\mu,\nu$ lassen sich aber auf unendlich viele Arten so bestimmen, dass $\mu m+\nu n=1$ gilt. Daraus ergeben sich immer neue Werte jenes Ausdrucks; untersuchen wir ihren Zusammenhang. Es gelte nicht nur $\mu m+\nu n=1$, sondern auch $\mu'm+\nu'n=1$, und wir setzen
\[
\mu(mb+nc)-\nu(ma+nb)=v,\quad
\mu'(mb+nc)-\nu'(ma+nb)=v'
\]
Multiplizieren wir die Gleichung $\mu m+\nu n=1$ mit $\mu'$, die andere Gleichung $\mu'm+\nu'n=1$ mit $\mu$ und subtrahieren, so ergibt sich $\mu'-\mu=n(\mu'\nu-\mu\nu')$. Multiplizieren wir ebenso die erste mit $\nu'$ und die zweite mit $\nu$, so ergibt die Subtraktion $\nu'-\nu=m(\mu\nu'-\mu'\nu)$. Daraus folgt sofort
\[
v'-v=(\mu'\nu-\mu\nu')(amm+2bmn+cnn)=(\mu'\nu-\mu\nu')M
\]
oder $v'\equiv v\ (\mathrm{mod.}\,M)$. Wie auch immer $\mu,\nu$ bestimmt werden, kann die Formel $\mu(mb+nc)-\nu(ma+nb)$ also keine \textit{verschiedenen}, das heißt inkongruenten, Werte des Ausdrucks $\surd(bb-ac)\ (\mathrm{mod.}\,M)$ liefern. Ist daher $v$ irgendein Wert dieser Formel, so sagen wir, dass die Darstellung der Zahl $M$ durch die Form $axx+2bxy+cyy$, bei der $x=m$, $y=n$ ist, \textit{zum Wert $v$ des Ausdrucks $\surd(bb-ac)(\mathit{mod.}\,M)$ gehört}. Außerdem lässt sich leicht zeigen: Ist ein Wert jener Formel $v$ und gilt $v'\equiv v\ (\mathrm{mod.}\,M)$, so kann man anstelle der Zahlen $\mu,\nu$, die $v$ liefern, andere Zahlen $\mu',\nu'$ wählen, die $v'$ liefern. Setzen wir nämlich
\[
\mu'=\mu+\frac{n(v'-v)}{M},\quad
\nu'=\nu-\frac{m(v'-v)}{M}
\]
so gilt
\[
\mu'm+\nu'n=\mu m+\nu n=1
\]
\SourceLeaf{128}
und der aus $\mu',\nu'$ hervorgehende Wert der Formel übertrifft den aus $\mu,\nu$ hervorgehenden um die Größe $(\mu'\nu-\mu\nu')M$. Diese wird $=(\mu m+\nu n)(v'-v)=v'-v$; für den erhaltenen Wert gilt also $=v'$.

\Article{156}
Haben wir zwei Darstellungen derselben Zahl $M$ durch dieselbe Form $(a,b,c)$, bei denen die Werte der Unbestimmten jeweils teilerfremd sind, so können sie entweder zu demselben Wert des Ausdrucks $\surd(bb-ac)\ (\mathrm{mod.}\,M)$ oder zu verschiedenen Werten gehören. Es sei
\[
M=amm+2bmn+cnn=am'm'+2bm'n'+cn'n'
\]
und
\[
\mu m+\nu n=1,\quad \mu'm'+\nu'n'=1
\]
Gilt
\[
\mu(mb+nc)-\nu(ma+nb)\equiv\mu'(m'b+n'c)-\nu'(m'a+n'b)\ (\mathrm{mod.}\,M)
\]
so bleibt die Kongruenz offenbar bestehen, welche anderen geeigneten Werte auch für $\mu,\nu$ und $\mu',\nu'$ gewählt werden. In diesem Fall sagen wir, dass beide Darstellungen zu \textit{demselben} Wert des Ausdrucks $\surd(bb-ac)\ (\mathrm{mod.}\,M)$ gehören. Gilt die Kongruenz dagegen bei irgendeiner Wahl von $\mu,\nu$ und $\mu',\nu'$ nicht, so gilt sie bei keiner, und die Darstellungen gehören zu \textit{verschiedenen} Werten. Ist hingegen
\[
\mu(mb+nc)-\nu(ma+nb)\equiv-(\mu'(m'b+n'c)-\nu'(m'a+n'b))
\]
so sagen wir, dass die Darstellungen zu \textit{entgegengesetzten} Werten des Ausdrucks $\surd(bb-ac)$ gehören. Alle diese Bezeichnungen verwenden wir auch dann, wenn mehrere Darstellungen derselben Zahl durch \textit{verschiedene} Formen mit derselben Determinante betrachtet werden.

\textit{Beispiel.} Gegeben sei die Form $(3,7,-8)$ mit der Determinante $=73$. Sie liefert folgende Darstellungen der Zahl $57$:
\[
3\mathbin{.}13^2+14\mathbin{.}13\mathbin{.}25-8\mathbin{.}25^2;\quad
3\mathbin{.}5^2+14\mathbin{.}5\mathbin{.}9-8\mathbin{.}9^2
\]
Für die erste können wir $\mu=2$, $\nu=-1$ setzen. Damit ergibt sich der Wert des Ausdrucks $\surd73\ (\mathrm{mod.}\,57)$, zu dem die Darstellung gehört:
\[
{}=2(13\mathbin{.}7-25\mathbin{.}8)+(13\mathbin{.}3+25\mathbin{.}7)=-4
\]
\SourceLeaf{129}
Ebenso findet man durch die Wahl $\mu=2$, $\nu=-1$, dass die zweite Darstellung zum Wert $+4$ gehört. Beide Darstellungen gehören also zu entgegengesetzten Werten.

Bevor wir weitergehen, bemerken wir, dass Formen mit der Determinante $=0$ aus den folgenden Untersuchungen vollständig ausgeschlossen sind. Sie würden nur die Geschlossenheit der Sätze stören und verlangen deshalb eine eigene Behandlung.

\begin{center}
\textit{Eine Form, die eine andere enthält oder in einer anderen enthalten ist; eigentliche und uneigentliche Transformation.}
\end{center}

\Article{157}
Lässt sich eine Form $F$ mit den Unbestimmten $x,y$ in eine andere Form $F'$ mit den Unbestimmten $x',y'$ durch Substitutionen
\[
x=\alpha x'+\beta y',\quad y=\gamma x'+\delta y'
\]
überführen, bei denen $\alpha,\beta,\gamma,\delta$ ganze Zahlen sind, so sagen wir, dass die erste die zweite \textit{enthält} oder dass die zweite \textit{in der ersten enthalten ist}. Die Form $F$ sei
\[
axx+2bxy+cyy
\]
und die Form $F'$ sei
\[
a'x'x'+2b'x'y'+c'y'y'
\]
Dann gelten die folgenden drei Gleichungen:
\[
\begin{aligned}
a'&=a\alpha\alpha+2b\alpha\gamma+c\gamma\gamma\\
b'&=a\alpha\beta+b(\alpha\delta+\beta\gamma)+c\gamma\delta\\
c'&=a\beta\beta+2b\beta\delta+c\delta\delta
\end{aligned}
\]
Multipliziert man die zweite Gleichung mit sich selbst und die erste mit der dritten und subtrahiert dann unter Weglassen der sich aufhebenden Glieder, so ergibt sich
\[
b'b'-a'c'=(bb-ac)(\alpha\delta-\beta\gamma)^2
\]
Daraus folgt, dass die Determinante der Form $F'$ durch die Determinante der Form $F$ teilbar und der Quotient eine Quadratzahl ist. Diese Determinanten haben also offenbar \textit{dasselbe Vorzeichen}. Lässt sich außerdem die Form $F'$ durch eine ähnliche Substitution in die Form $F$ überführen, ist also sowohl $F'$ in $F$ als auch $F$ in $F'$ enthalten, so sind
\SourceLeaf{130}
die Determinanten der Formen gleich\HistoricalNote{*)}{Aus der vorangehenden Rechnung geht klar hervor, dass diese Aussage auch für Formen mit der Determinante $=0$ gilt. Die Gleichung $(\alpha\delta-\beta\gamma)^2=1$ darf aber nicht auf diesen Fall ausgedehnt werden.} und es gilt $(\alpha\delta-\beta\gamma)^2=1$. In diesem Fall nennen wir die Formen \textit{äquivalent}. Die Gleichheit der Determinanten ist also eine notwendige Bedingung für die Äquivalenz von Formen, obwohl diese daraus allein keineswegs folgt. --- Die Substitution $x=\alpha x'+\beta y'$, $y=\gamma x'+\delta y'$ nennen wir eine \textit{eigentliche Transformation}, wenn $\alpha\delta-\beta\gamma$ positiv ist, und eine \textit{uneigentliche}, wenn $\alpha\delta-\beta\gamma$ negativ ist. Wir sagen, dass die Form $F'$ \textit{eigentlich} oder \textit{uneigentlich} in der Form $F$ enthalten ist, wenn $F$ durch eine eigentliche oder uneigentliche Transformation in die Form $F'$ überführt werden kann. Sind also die Formen $F,F'$ äquivalent, so gilt $(\alpha\delta-\beta\gamma)^2=1$ und daher bei einer eigentlichen Transformation $\alpha\delta-\beta\gamma=+1$, bei einer uneigentlichen $=-1$. --- Sind mehrere Transformationen sämtlich eigentlich oder sämtlich uneigentlich, so nennen wir sie \textit{gleichartig}; eine eigentliche und eine uneigentliche dagegen \textit{ungleichartig}.

\begin{center}
\textit{Eigentliche und uneigentliche Äquivalenz.}
\end{center}

\Article{158}
\textit{Haben die Formen $F,F'$ gleiche Determinanten und ist $F'$ in $F$ enthalten, so ist auch $F$ in $F'$ enthalten, und zwar eigentlich oder uneigentlich, je nachdem, ob $F'$ eigentlich oder uneigentlich in $F$ enthalten ist.} Die Form $F$ gehe in $F'$ über, indem wir setzen:
\[
x=\alpha x'+\beta y',\quad y=\gamma x'+\delta y'
\]
Dann geht $F'$ in $F$ über, wenn man
\[
x'=\delta x-\beta y,\quad y'=-\gamma x+\alpha y
\]
setzt. Denn diese Substitution liefert aus $F'$ offenbar dasselbe, was aus $F$ durch die Substitution
\[
x=\alpha(\delta x-\beta y)+\beta(-\gamma x+\alpha y),\quad
y=\gamma(\delta x-\beta y)+\delta(-\gamma x+\alpha y)
\]
oder
\[
x=(\alpha\delta-\beta\gamma)x,\quad y=(\alpha\delta-\beta\gamma)y
\]
entsteht. Dadurch wird $F$ offenbar in $(\alpha\delta-\beta\gamma)^2F$, das heißt wieder in $F$ überführt (voriger Artikel). Die zweite Transformation ist offenbar eigentlich oder uneigentlich, je nachdem, ob die erste eigentlich oder uneigentlich ist.

Sind sowohl $F'$ in $F$ als auch $F$ in $F'$ \textit{eigentlich} enthalten, so nennen wir die Formen \textit{eigentlich \SourceLeaf{131}äquivalent}; sind sie jeweils uneigentlich ineinander enthalten, so nennen wir sie \textit{uneigentlich äquivalent}. --- Der Nutzen dieser Unterscheidungen wird sich bald zeigen.

\Needspace{6\baselineskip}
\textit{Beispiel.} Die Form $2xx-8xy+3yy$ geht durch die Substitutionen $x=2x'+y'$, $y=3x'+2y'$ in die Form $-13x'x'-12x'y'-2y'y'$ über, diese wiederum in jene durch die Substitutionen $x'=2x-y$, $y'=-3x+2y$. Die Formen $(2,-4,3)$, $(-13,-6,-2)$ sind also \textit{eigentlich äquivalent}.

Die Probleme, denen wir uns nun zuwenden, sind folgende. I. Für zwei beliebige gegebene Formen mit derselben Determinante ist zu untersuchen, ob sie äquivalent sind oder nicht, und ob eigentlich, uneigentlich oder auf beide Arten, denn auch dies ist möglich. Haben sie verschiedene Determinanten, so ist zu untersuchen, ob wenigstens eine die andere enthält, eigentlich, uneigentlich oder auf beide Arten. Schließlich sind alle Transformationen der einen in die andere zu finden, sowohl die eigentlichen als auch die uneigentlichen. II. Für eine beliebige gegebene Form ist zu ermitteln, ob eine gegebene Zahl durch sie dargestellt werden kann, und es sind alle Darstellungen anzugeben. Da Formen mit negativer Determinante hier aber eine andere Methode verlangen als Formen mit positiver Determinante, werden wir zuerst das behandeln, was beiden gemeinsam ist, und danach die Formen jeder Art gesondert betrachten.
